\section{Knowledge Distillation: Foundations} \label{sec:23-background-kdfoundations}

\subsection{Seminal Work: Hinton et al.\ (2015)}
\label{sec:23-background-kdfoundations-hinton}

\begin{foldable}
  Knowledge distillation~\cite{hinton2015distilling} transfers knowledge from a large, pre-trained \emph{teacher} network $T$ to a smaller \emph{student} network $S$ by using the teacher's output distribution as a training target, in addition to or in place of hard ground-truth labels.
  The central insight is that a teacher's softmax output over all $K$ classes encodes substantially more information than the corresponding one-hot label.
  When the teacher evaluates an image of the digit ``8'', it assigns non-trivial probability to ``3'' and near-zero probability to ``1'', implicitly encoding the fact that eights and threes are more perceptually similar than eights and ones.
  Hinton et al.\ termed this inter-class relational information \emph{dark knowledge}: it is present in every teacher query but is completely discarded when supervision is collapsed to a single correct-class label.
\end{foldable}

\begin{foldable}
  To amplify dark knowledge during training, a temperature parameter $\tau \geq 1$ is applied to the teacher's pre-softmax logits before computing the output distribution:
  \begin{equation}
    q_k^{(\tau)} = \frac{\exp(z_k^{T} / \tau)}{\sum_{j=1}^{K} \exp(z_j^{T} / \tau)},
    \label{eq:softmax-temperature}
  \end{equation}
  where $z_k^{T}$ is the teacher logit for class $k$.
  Higher temperatures produce softer distributions, assigning greater relative weight to lower-probability classes and making the similarity structure more visible to the student.
  The student's logits are passed through the same temperature to produce $p_k^{(\tau)}$, and the student minimises the KL divergence between the two softened distributions.
\end{foldable}

\begin{foldable}
  The full training objective combines soft and hard supervision:
  \begin{equation}
    \mathcal{L}_{\mathrm{KD}}
    = \alpha \, \tau^{2} \, \mathrm{KL}\!\left(\mathbf{q}^{(\tau)} \,\|\, \mathbf{p}^{(\tau)}\right)
    + (1 - \alpha) \, \mathcal{L}_{\mathrm{CE}}\!\left(\mathbf{p}^{(1)}, y\right),
    \label{eq:kd-loss}
  \end{equation}
  where $\mathcal{L}_{\mathrm{CE}}$ is the standard cross-entropy loss against the hard label $y$, and $\alpha \in [0,1]$ balances the two terms.
  The $\tau^{2}$ scaling compensates for the reduction in gradient magnitude caused by temperature-softening~\cite{hinton2015distilling}.
  In practice, $\tau$ is set between 2 and 10 and $\alpha$ is tuned per task, with larger $\alpha$ giving more weight to the soft signal.
\end{foldable}

\begin{foldable}
  The empirical gain from KD is well established: students trained with soft targets consistently outperform counterparts trained on hard labels alone, even when the student architecture is identical to the teacher~\cite{hinton2015distilling}.
  This \emph{self-distillation} result confirms that the benefit is not attributable to capacity mismatch but to the information content of the soft target itself.
\end{foldable}

\subsection{Taxonomy of Knowledge Distillation Methods}
\label{sec:23-background-kdfoundations-taxonomy}

\begin{foldable}
  The field has expanded well beyond the original logits-based formulation, with methods differing primarily in the \emph{type of teacher signal} they transfer.
  Three broad families have emerged.
\end{foldable}

\begin{foldable}
  \textbf{Logits-based distillation} follows Hinton et al.\ directly: the student learns to match the teacher's output class distribution.
  The soft-target loss in Equation~\ref{eq:kd-loss} is the canonical form.
  Variants include using the full un-softened logit vector as a regression target (removing the temperature hyperparameter at the cost of sensitivity to logit scale) and born-again networks, in which a student of the same architecture as the teacher is sequentially re-distilled from the previous generation, yielding ensembles that outperform the original teacher~\cite{ba2014deep}.
  Logits-based methods require only the teacher's output probability vector and are therefore compatible with any setting in which those outputs are accessible.
\end{foldable}

\begin{foldable}
  \textbf{Feature-based distillation} transfers knowledge from the teacher's intermediate representations, typically by adding an auxiliary loss that minimises the distance between aligned teacher and student feature maps.
  FitNets~\cite{romero2014fitnets} introduced the idea of matching intermediate hidden layers using a thin regressor, enabling the training of networks deeper and thinner than the teacher.
  Attention Transfer~\cite{zagoruyko2016paying} distils spatial attention maps derived from convolutional activations, capturing \emph{where} the teacher focuses rather than \emph{what} it predicts.
  PKT~\cite{passalis2018learning} reframes feature transfer as probabilistic knowledge: it matches the kernel density estimates of teacher and student feature distributions, encoding the distributional geometry of the representation space rather than individual activation values.
  Feature-based methods consistently improve on logits-only baselines on standard benchmarks, but they require access to the teacher's internal activations --- a requirement that is violated whenever the teacher is deployed as an opaque API.
\end{foldable}

\begin{foldable}
  The difficulty this poses goes deeper than a mere engineering inconvenience.
  Feature-based distillation aligns teacher and student in representation space, which presupposes that the two spaces are compatible: either they share the same dimensionality, or a learned projector can be trained to bridge them.
  Designing such a projector requires knowing the teacher's activation dimensionality and, ideally, the semantic structure of its representation space --- which regions cluster together, which directions encode discriminative axes, and how the geometry evolves across layers.
  When the teacher is an opaque \ac{API}, none of this information is available. The practitioner cannot inspect the layer topology, cannot determine the width of any intermediate feature map, and cannot characterise the manifold structure of the representation space.
  The projector design problem is therefore not just unconstrained but fundamentally ill-posed: there is no signal with which to train a bridge from student to teacher space, because the target space is entirely unobservable.
  FitNets~\cite{romero2014fitnets} illustrated the potential of such projection-based alignment in the white-box setting, and AT~\cite{zagoruyko2016paying} demonstrated that even compressed, spatially-pooled representations carry transferable information --- yet both methods rely on a clearly defined teacher layer from which to extract the supervision signal.
  Remove that access and the entire apparatus collapses: there is no anchor in the teacher's representation space to which the student can be steered.
  This is a fundamental incompatibility between the feature-based paradigm and the restricted-access settings considered in \S\ref{sec:24-background-kdrestricted}, not a problem that better optimization or richer student capacity can resolve.
\end{foldable}

\begin{foldable}
  \textbf{Relation-based distillation} transfers higher-order structure: pairwise or group-wise relationships among data samples in the teacher's representation space.
  RKD~\cite{park2019relational} penalises discrepancies in the distance and angle relationships between pairs of training examples in teacher and student embedding spaces, arguing that relational information is more transferable across architectures than pointwise feature values.
  Self-supervised knowledge distillation methods extend this idea by using contrastive objectives to enforce that the student's embedding space preserves the teacher's instance discrimination structure.
  Relation-based methods also require at least intermediate-level teacher access to compute the pairwise statistics.
\end{foldable}

\begin{foldable}
  Relation-based methods carry an additional failure mode that is structural rather than purely an access problem: the relational signal is inherently batch-dependent, and its quality degrades sharply under the conditions that characterise restricted-access settings.
  RKD~\cite{park2019relational} constructs supervision from pairwise distance and angle relationships among samples within a minibatch, which means the number of meaningful training signals scales as $O(B^2)$ in batch size $B$.
  In few-shot or black-box settings, however, the budget of queries to the teacher is severely constrained, forcing small batches with sparse class coverage.
  A minibatch drawn from only a handful of classes cannot faithfully approximate the relational geometry of the full data distribution: the pairwise statistics it yields are noisy estimates of an unknown population geometry, and the student is supervised toward a distorted relational target.
  Even if batch size were not a constraint, a subtler incompatibility remains.
  Relation-based methods implicitly assume that teacher and student share a coherent notion of proximity in embedding space: that samples which are close under the teacher's metric are also the samples the student should treat as close.
  But when the teacher is inaccessible, the student's embedding space is shaped exclusively by hard-label cross-entropy supervision, which enforces correct classification but imposes no constraint on the metric structure of the learned representations.
  The student's geometry and the teacher's geometry may be entirely incommensurable --- not merely rotated or scaled, but organised along different discriminative axes.
  Under those conditions, pairwise relational supervision provides a target that the student cannot meaningfully pursue, because the spaces in which the two sets of pairwise relationships are defined bear no systematic correspondence to one another.
  Together, these scalability and geometry-mismatch problems explain why the restricted-access settings surveyed in \S\ref{sec:24-background-kdrestricted} require methods that abandon the relational transfer paradigm and instead seek signals that are both extractable from opaque outputs and informative enough to guide compact student training.
\end{foldable}

\begin{table}[t]
  \centering
  \caption{Summary of the three KD families by signal type and access requirement.}
  \label{tab:kd-taxonomy}
  \begin{tabular}{llll}
    \hline
    \textbf{Family} & \textbf{Signal transferred} & \textbf{Teacher access} & \textbf{Data required} \\
    \hline
    Logits-based    & Output class distribution   & Outputs only            & Training set $\mathcal{D}$ \\
    Feature-based   & Intermediate activations    & Internal layers         & Training set $\mathcal{D}$ \\
    Relation-based  & Pairwise embedding geometry & Internal layers         & Training set $\mathcal{D}$ \\
    \hline
  \end{tabular}
\end{table}

\begin{foldable}
  Table~\ref{tab:kd-taxonomy} summarises the three families.
  All standard KD methods share two implicit assumptions: the full training dataset $\mathcal{D}$ is available to construct transfer queries, and the teacher exposes at least its output probabilities.
  The next subsection examines what happens when these assumptions are violated.
\end{foldable}

\subsection{Standard Assumptions and Their Limitations}
\label{sec:23-background-kdfoundations-assumptions}

\begin{foldable}
  Standard KD rests on two foundational assumptions that are frequently violated in practice.
\end{foldable}

\begin{foldable}
  \textbf{Assumption 1: the full training dataset is available.}
  KD methods use training examples to query the teacher, collect soft targets, and train the student.
  In many realistic settings, the original training data is inaccessible: it may contain private user records (medical imaging, financial transactions), it may be proprietary to the model provider, or it may simply have been discarded after training to reduce storage costs.
  Data regulations such as GDPR further restrict the retention and transfer of training sets across organizational boundaries.
  When $\mathcal{D}$ is unavailable, the student must be trained on some surrogate: a small held-out subset, synthetic data generated by a generative model, or public proxy data from a related domain.
\end{foldable}

\begin{foldable}
  \textbf{Assumption 2: the teacher is accessible at the logit (or deeper) level.}
  All three KD families in Table~\ref{tab:kd-taxonomy} require at minimum the teacher's output probability vector over all $K$ classes.
  Feature-based and relation-based methods additionally require access to internal activations.
  Modern production models are routinely deployed behind APIs that return only a top-1 predicted label, or the top-3 to top-5 labels with associated confidence scores, to protect proprietary model internals and limit adversarial probing.
  In this setting, neither soft targets nor intermediate features are accessible, and the information bottleneck is severe: a single hard label provides $\log_2 K$ bits of information per query, compared to a soft target vector which, for a well-calibrated teacher at moderate temperature, encodes rich inter-class similarity structure.
\end{foldable}

\begin{foldable}
  The combination of these two failures defines the landscape of practical KD: feature-based and relation-based methods are ruled out when teacher internals are inaccessible;
  logits-only methods are ruled out when the API returns only hard labels; and all methods degrade when training data is replaced by a small or non-representative surrogate.
  \S\ref{sec:24-background-kdrestricted} surveys the methods that have been developed for each restricted-access setting and identifies the open problems that this thesis addresses.
\end{foldable}

\begin{foldable}
  It is also worth noting here that KD addresses \emph{model} compression: the goal is to produce a smaller model that approximates the teacher's predictive behaviour.
  A complementary problem, \emph{dataset} compression, asks instead how to select or synthesise a compact training corpus that, when used to train a fresh model, approximates the performance of training on the full dataset.
  These two problems are related but distinct; their relationship and the specific challenges of dataset compression in the text domain are taken up in \S\ref{sec:26-background-datasetdistillation}.
\end{foldable}
